\section{The actual vacuum identity and algebraic germs}
\label{sec:vacuum-identity}

This section supplies the vacuum and four-copy inputs retained in
TL18. Fix the literal arithmetic datum \(D\), a curve ideal
\(\mathfrak c\), its prescribed discriminant root \(\delta\), and
an actual \(\mathscr B:\texttt{Base}(D,\mathfrak c,\delta)\).
We assume the source's \(\texttt{Crv}\) conditions. Every period,
half-point, isometry and correction root below is the one in
\(\mathscr B\); no replacement base is chosen.

Write \(K=\mathbb Q(\sqrt{-d})\subset E\), \(A=D.\mathrm{host}.\mathrm{AA0}\),
\(n=|\mathrm{Tau}|\), \(\iota:K\hookrightarrow\mathbb C\), and
\(\operatorname{emb}(b)=(\tau(b))_\tau\). Thus
\(\vartheta=(1+i\sqrt d)/2\), \(d\equiv7\pmod {16}\), and
\[
 \Lambda=\delta\iota(\mathfrak c),\qquad
 a_\Lambda=\operatorname{area}(\Lambda),\qquad
 H(z,w)=a_\Lambda^{-1}\sum_\tau z_\tau\overline{w_\tau},
 \qquad E_H(z,w)=\operatorname{Im}H(z,w).
 \tag{VI.1}
\]
The form is linear in its first variable. The supplied half-point is
\(t=\omega_0/2\); it satisfies the three actual unkilled clauses.
The arithmetic host, including its trace-dual fractional ideal, remains
an input. The one-variable analytic and CM assertions used below are
the reconstructed PA and TA assertions, including TA14.

\paragraph{VI1: the same arithmetic lattices, complex map, and correction.}
The additive subgroup called \(\texttt{lamM}\) is exactly
\[
             \mathcal M=\delta\operatorname{emb}(\mathfrak c A).
 \tag{VI1.1}
\]
Indeed its defining generators are
\(\delta\iota(c)\operatorname{emb}(b)=
\delta\operatorname{emb}(cb)\), with \(c\in\mathfrak c\),
\(b\in A\); their finite sums are the ideal product
\(\mathfrak c A\). The closure in this definition is additive
subgroup generation. It is not a topological closure.

Let \(\gamma_0\) and \(g=\mathscr B.g_{\mathbb C}\) be the supplied
maps. PA4 proves that the vectors
\(\operatorname{emb}_4(x)=(\tau(x_a))_{a,\tau}\) span
\(\mathbb C^{4n}\), and that the unique complex extension of
\(\gamma_0\) has an algebraic matrix and is an isometry for the
sum Hermitian form. Since the supplied \(g\) has exactly that
reading on these spanning vectors, it is this same map. In particular
both \(g\) and \(g^{-1}\) preserve algebraic coordinate vectors.

Put
\[
 \mathcal S=\Lambda^{4n},\qquad
 \mathcal L=g(\mathcal M^4),\qquad
 L_K=\gamma_0(A^4)\subset K^{4n}.
 \tag{VI1.2}
\]
The same reading and \(K\)-linearity give
\(\mathcal L=\delta\iota(\mathfrak c L_K)\), coordinate by coordinate.
The denominator argument of PA10 gives an integer \(N>0\) with
\(NL_K\subset\mathcal O_K^{4n}\) and
\(N\mathcal O_K^{4n}\subset L_K\). Multiplying these inclusions
by \(\mathfrak c\) and applying \(\delta\iota\) proves
\[
                N\mathcal L\subset\mathcal S,\qquad
                N\mathcal S\subset\mathcal L.
 \tag{VI1.3}
\]
Both are full, finitely generated lattices, and both relative
quotients by their intersection are finite. Thus the source
indices \(\texttt{idx}_1,\texttt{idx}_2\) are positive. The first
inclusion also proves that each coordinate of each \(y\in\mathcal L\)
is a rational-lattice point of \(\Lambda\).

Finally PA10 applies to the actual \(\mathscr B.kc\). Its determinant
exponents vanish outside finitely many rational primes, and its
chosen roots satisfy their nonzero algebraic \(2h_K\)-power
equations. Hence the actual correction is
\[
                  \kappa=\prod_p r_p^{-y_p}
                         \in\overline{\mathbb Q}^{\,\times}.
 \tag{VI1.4}
\]
This is a genuinely finite product on the prescribed prime subtype.
It retains the chosen roots, without asserting that different
independent root choices give the same correction.

\paragraph{VI2: the full period lattices are self-dual with the stated multiplier.}
Here is the integral calculation behind the self-duality input.
Let \(L\subset E^N\), \(L'\subset E^{N'}\) be full self-dual
\(\mathcal O_E\)-lattices and put
\[
 P=\operatorname{Hom}_{\mathcal O_E}(L,L')
       \subset M_{N',N}(E).
 \tag{VI2.1}
\]
The pairing on matrices is \(h(T,U)=\operatorname{tr}(TU^*)\).
On their complex embeddings, (VI.1) is extended by summing
over the row and column indices as well as \(\tau\).
Its \(\mathcal O_E\)-dual on \(P\) is \(P\) itself. To check
this, localize at each finite prime. Choose bases for the two
modules and, where conjugation changes the prime, for the modules
at the conjugate prime. Self-duality identifies the corresponding
dual bases. The matrix units are then dual for the trace pairing.
This argument requires perfect integral pairings, not orthonormal
integral bases.

The host's trace-dual identity for \(A\) is equivalently
\(A\bar A=\mathfrak D_{E/K}^{-1}\). The trace-dual ideal formula,
applied to the preceding matrix units, therefore gives
\[
 (AP)^{\#_K}=AP,\qquad
 (\mathfrak c AP)^{\#_K}=\bar{\mathfrak c}^{-1}AP,\qquad
 (\operatorname N\mathfrak c)\,
       (\mathfrak c AP)^{\#_K}=\mathfrak c AP.
 \tag{VI2.2}
\]
Here \(\#_K\) tests \(\operatorname{Tr}_{E/K}h\) for
\(\mathcal O_K\)-valuedness. The last equality uses the actual
quadratic ideal identity
\(\mathfrak c\bar{\mathfrak c}=
(\operatorname N\mathfrak c)\mathcal O_K\).

The source lattice is precisely
\[
 \mathcal A(L,L')=\texttt{La}(L,L')
                 =\delta\operatorname{emb}(\mathfrak c AP).
 \tag{VI2.3}
\]
Expanding its defining generators proves this equality in both
directions. PA1 gives
\(a_\Lambda=|\delta|^2(\operatorname N\mathfrak c)\sqrt d/2\).
For \(b,c\) in the rational matrix space, the resulting alternating
pairing is
\[
 E_H(\delta\operatorname{emb}b,\delta\operatorname{emb}c)
 =\frac{1}{\operatorname N\mathfrak c}
    \operatorname{Tr}_{K/\mathbb Q}
       \left(\frac{\operatorname{Tr}_{E/K}h(b,c)}
                         {\sqrt{-d}}\right).
 \tag{VI2.4}
\]
The codifferent of \(K/\mathbb Q\) is
\((\sqrt{-d})^{-1}\mathcal O_K\). Testing (VI2.4) against all
\(\mathcal O_K\)-multiples of \(c\), and using (VI2.2), proves
that the alternating integral dual of \(\mathcal A(L,L')\)
is itself. This test also handles a vector initially given in
the full real space: in an integral lattice basis its pairings
are integral and its coordinates are obtained by inverting a
nonsingular rational pairing matrix. It therefore lies in the
rational span, where (VI2.2) applies.

These are full lattices. Any two of the modules in (VI2.3) have
the same rational span and admit common integer denominator
bounds, so the resulting period lattices are commensurable.
This proves the self-duality and commensurability clauses of the
actual \(\texttt{LatAll}\). In particular \(\mathcal M\) is
self-dual for \(E_H\); its fourth power is self-dual, and the
isometry \(g\) makes \(\mathcal L\) self-dual. The coordinate
period lattice \(\mathcal S\) is self-dual as well: on a period
basis the alternating matrix has off-diagonal entries \(1,-1\),
up to orientation, and one tests coordinates separately.

Every lattice just described is stable under \(\vartheta\).
For its vectors \(x,y\), write
\(e=E_H(x,y)\in\mathbb Z\) and
\(k=E_H(\vartheta x,y)\in\mathbb Z\). Since
\[
 \sqrt d\,\operatorname{Re}H(x,y)=2k-e,
 \tag{VI2.5}
\]
the quadratic expression
\(q(x)=\sqrt d\,\operatorname{Re}H(x,x)/2\) is an integer,
and
\[
 q(x+y)-q(x)-q(y)-E_H(x,y)\in2\mathbb Z.
 \tag{VI2.6}
\]
Consequently \(\alpha(x)=\exp(\pi i q(x))\) takes values in
\(\{1,-1\}\) and is exactly the required semicharacter:
\[
 \alpha(x+y)=\alpha(x)\alpha(y)\exp(\pi iE_H(x,y)).
 \tag{VI2.7}
\]
On every actual \(\texttt{La}\), the source's
\((-1)^{\operatorname{round}(q(x))}\) equals this \(\alpha(x)\),
because \(q(x)\) is integral. Thus this computation supplies the
literal multiplier, including its rounding convention.

\paragraph{VI3: the chosen half-point supplies the standard character and StdOK.}
For the one-dimensional period lattice let
\[
 \eta(w)=
 \begin{cases}1&w\in2\Lambda,\\-1&w\notin2\Lambda,\end{cases}
 \qquad
 \chi_0(w)=\eta(w)\exp(\pi iE_H(w,\omega_0)).
 \tag{VI3.1}
\]
The unkilled conditions imply that
\(\omega_0,\vartheta\omega_0,\omega_0+\vartheta\omega_0\)
are all nonzero in \(\Lambda/2\Lambda\). Thus the first two
form a basis of this two-dimensional \(\mathbb F_2\)-space.
Its alternating form is nondegenerate, so
\[
 q(\omega_0)=E_H(\vartheta\omega_0,\omega_0)
                         \quad\hbox{is odd}.
 \tag{VI3.2}
\]
Also \(q(\vartheta\omega_0)=
\frac{1+d}{4}q(\omega_0)\) is even, since \(d\equiv7\pmod8\).
It follows directly from (VI3.1) that \(\alpha\) and \(\chi_0\)
have the same values on both basis vectors: respectively \(-1\)
and \(1\). Their product is an ordinary character trivial on
\(2\Lambda\), because their semicharacter factors cancel.
They therefore agree on all of \(\Lambda\).
Taking products over the coordinates gives
\[
       \alpha|_{\mathcal S}=\texttt{stdChar}(L,\omega_0).
 \tag{VI3.3}
\]

We can now verify all ten fields of \(\texttt{StdData}\), hence
the actual \(\mathscr B.\texttt{StdOK}\).
The period and nonzero half-period are supplied by the actual
unkilled point. The isometry is VI1; \(\alpha(0)=1\);
the standard character is (VI3.3); and VI2 proves the image
lattice's semicharacter and self-duality. For every four-block
vector, isometry and the block sum give the exact identity
\[
             \alpha(gz)=\prod_{a=0}^3\alpha(z_a).
 \tag{VI3.4}
\]
The two positive relative indices are VI1. These verifications
prove \(\texttt{stdOK}\) and \(\texttt{stdOK\_of\_crv}\);
they do not leave \(\texttt{StdOK}\) as an independent premise.

\paragraph{VI4: the Heisenberg formula and a well-defined finite average.}
For a Hermitian form linear in its first variable define
\[
 U_xF(z)=F(z+x)
          \exp\!\left(-\pi H(z,x)-\frac{\pi}{2}H(x,x)\right).
 \tag{VI4.1}
\]
Expanding the exponent, with
\(H(x,y)-H(y,x)=2iE_H(x,y)\), gives
\[
 U_xU_y=e^{-\pi iE_H(x,y)}U_{x+y},\qquad
 U_xU_y=e^{-2\pi iE_H(x,y)}U_yU_x.
 \tag{VI4.2}
\]
For a lattice with semicharacter \(\alpha\), put
\(B_x=\alpha(x)U_x\). Equations (VI2.6)--(VI2.7) and integrality
of \(E_H\) show \(B_xB_y=B_{x+y}\) for vectors of that lattice.
In particular the theta condition is precisely \(B_xF=F\).

Let \(\Lambda_1,\Lambda_2\) be commensurable lattices carrying
the same semicharacter, and let \(F\) satisfy the theta condition
for \(\Lambda_1\). If \(y\in\Lambda_2\) and
\(d\in\Lambda_1\cap\Lambda_2\), then
\(B_{y+d}F=B_yB_dF=B_yF\). Thus
\[
 I_{12}F=
 \frac1{\sqrt{r_{12}}}
 \sum_{[y]\in\Lambda_2/(\Lambda_1\cap\Lambda_2)}B_yF,
 \qquad r_{12}=[\Lambda_2:\Lambda_1\cap\Lambda_2],
 \tag{VI4.3}
\]
is independent of every representative choice. The square root
is the positive real square root. Translation by any
\(x\in\Lambda_2\) permutes the summands, proving
\(I_{12}F\) is fixed by \(\Lambda_2\). Every summand is entire
when \(F\) is entire, so the finite sum is entire.
This is exactly the source's \(\texttt{inter}\), whose
\texttt{finsum} is an ordinary finite sum here.

\paragraph{VI5: finite Fourier inversion, nonvanishing, and the positive normalization.}
Suppose both lattices in VI4 are self-dual. The pairing
\[
 (\Lambda_1/(\Lambda_1\cap\Lambda_2))\times
 (\Lambda_2/(\Lambda_1\cap\Lambda_2))
 \longrightarrow\mathbb C^\times,\qquad
 (x,y)\longmapsto e^{-2\pi iE_H(x,y)}
 \tag{VI5.1}
\]
is well-defined and nondegenerate. For example, if its character
is trivial for every \(x\in\Lambda_1\), then
\(y\in\Lambda_1^\#=\Lambda_1\), so its quotient class is zero.
The other variable is identical.

A nontrivial character on a finite group has sum zero: translate
the sum by a point where its value is not \(1\). A trivial
character has sum equal to the group order. Summing (VI5.1)
in both orders consequently proves
\[
                   r_{12}=r_{21}=:r.
 \tag{VI5.2}
\]
Now expand \(I_{21}I_{12}F\). Move \(B_x\), \(x\in\Lambda_1\),
past \(B_y\) by (VI4.2), and use \(B_xF=F\).
The sum over \(x\) vanishes unless \(y\) is the zero quotient
class, where it equals \(r\). The two normalization factors
are \(r^{-1/2}\), so
\[
                         I_{21}I_{12}F=F.
 \tag{VI5.3}
\]
This proves both inversion and nonvanishing on a nonzero
theta vector. It also transports a one-dimensional theta
space to a one-dimensional theta space, in both directions.

For precision, the same positive normalization preserves the
Gaussian-weighted torus mean. On the common finite cover by
\(\Lambda_1\cap\Lambda_2\), use the normalized Haar mean of
\(F\overline G e^{-\pi H(z,z)}\). Directly from (VI4.1),
\[
 |U_xF(z)|^2e^{-\pi H(z,z)}
                 =|F(z+x)|^2e^{-\pi H(z+x,z+x)}.
 \tag{VI5.4}
\]
Thus \(B_x\) is unitary for this mean. The translates \(B_yF\)
belong to distinct character spaces for
\(\Lambda_1/(\Lambda_1\cap\Lambda_2)\), by (VI5.1), and hence
are pairwise orthogonal: applying an element with distinct
eigenvalues to their inner product proves it is zero.
All their squared norms equal that of \(F\). Equation (VI4.3)
therefore preserves the mean squared norm. The normalized mean
of a function periodic for either original lattice is unchanged
on this finite cover. This proves the stated isometry without
introducing an unspecified Fourier scalar or an analytic
nonvanishing assumption. The source's \(\texttt{normSq}\) is
exactly this normalized average over its chosen fundamental
domain. Fullness provides a bounded fundamental domain with
positive finite volume; periodicity makes the average independent
of that choice. Thus the argument gives the literal
\(\texttt{IopIso}\) norm equality.

\paragraph{VI6: the standard product theta space is exactly one line.}
Let \(s_\Lambda\) be the actual normalized odd theta function and
put
\[
 \theta_0(z)=\frac{U_t s_\Lambda(z)}{(U_t s_\Lambda)(0)},
 \qquad t=\omega_0/2.
 \tag{VI6.1}
\]
TA14 and PA15 give the source's function, its entirety and
simple zeros at \(-t+\Lambda\). The denominator is nonzero
because \(t\notin\Lambda\); \(\theta_0(0)=1\). Its multiplier
is precisely (VI3.1). This also identifies the period-pair
and \(\texttt{GL3Theta.th0}\) definitions on the same lattice.
TA12 keeps this normalized function on the prescribed
distinguished half-point class under coefficient conjugacy.
Thus the arithmetic input used later is attached to this
function, rather than to an unspecified generator of its line.

An entire one-variable function \(F\) with this multiplier
is a scalar multiple of \(\theta_0\). Indeed \(F/\theta_0\)
is elliptic and has at most one possible simple pole on the
torus. The sum of residues on a fundamental parallelogram is
zero, by cancellation of opposite edges. That sole possible
residue is therefore zero, so the pole is removable.
The quotient is entire and periodic, hence constant.
Its value at zero is \(F(0)\).

For finitely many coordinates put
\(F_0(z)=\prod_j\theta_0(z_j)\). The exponential in (VI4.1)
splits as a product, so \(F_0\) has the required standard
multiplier and \(F_0(0)=1\). Conversely fix all but one
coordinate of a standard theta function. Its remaining
one-variable slice has the preceding multiplier, so replacing
that coordinate by zero extracts the factor \(\theta_0\).
Repeating over all coordinates proves
\[
               F(z)=F(0)\prod_j\theta_0(z_j).
 \tag{VI6.2}
\]
This proves the actual standard theta-line theorem, together
with its normalized nonzero generator.

\paragraph{VI7: constructing the four-copy root and the actual chosen vacuum.}
Use the lattices and maps in VI1 and the multiplier in VI3.
By VI5--VI6 the function
\[
 K_4(z)=
 \frac1{\sqrt r}
 \sum_{[y]\in\mathcal L/(\mathcal L\cap\mathcal S)}
       \alpha(y)\,(U_yF_0)(gz)
 \tag{VI7.1}
\]
is nonzero. It spans the theta space for \(\mathcal M^4\),
the four-block form, and the product multiplier; the exact
change of variables is \(g\), by (VI3.4).
Set \(P=\kappa K_4\). It too is a nonzero generator.

Choose a point \(z^*=(z_0^*,z_1^*,z_2^*,z_3^*)\) with
\(P(z^*)\ne0\), and define
\(f_1(v)=P(v,z_1^*,z_2^*,z_3^*)\).
The block form and \(\alpha(0)=1\) show directly that \(f_1\)
is in the one-copy theta space. It is nonzero at \(z_0^*\).
Its tensor product \(f_1^{\otimes4}\) is in the four-copy line,
so \(f_1^{\otimes4}=cP\). Here \(c\ne0\): evaluate the left
side with all four variables equal to \(z_0^*\).
Choose \(b\in\mathbb C\) with \(b^4=c^{-1}\). Then
\(f=bf_1\) satisfies
\[
                   f^{\otimes4}=\kappa K_4.
 \tag{VI7.2}
\]
This takes a fourth root of one nonzero complex scalar; no
algebraicity over \(\mathbb Q\) is asserted for that scalar
or its chosen root, and no branch of a pointwise analytic
fourth root has been assumed.

The one-copy theta space is also a line. Given another vector
\(F\) in it, the mixed tensor \(F\otimes f\otimes f\otimes f\)
lies in the four-copy line. Evaluate its last three variables
at a point where \(f\ne0\) to obtain \(F=cf\).
If two such vectors satisfy (VI7.2), their ratio consequently
has fourth power \(1\).

We have proved existence, nonvanishing, spanning and the exact
normalization for every admissible choice. The literal
\(\mathscr B.\texttt{thMM}\) is the choice from this nonempty
predicate, so it satisfies (VI7.2) and is nonzero.
This proves both occurrences of \(\texttt{thMM\_spec}\), as
well as \(\texttt{thMM\_ne\_zero}\), with no independent
vacuum-existence premise.

\paragraph{VI8: the finite four-copy expression is an algebraic germ.}
Write \(f=\mathscr B.\texttt{thMM}\). For the coordinate form,
direct expansion of (VI4.1) gives
\[
                 (U_yF_0)(w)
                    =\prod_j (U_{y_j}\theta_0)(w_j).
 \tag{VI8.1}
\]
Every \(y_j\) appearing in (VI7.1) is a rational-lattice
point by (VI1.3). TA10, PA15 and the same half-point
normalization therefore give all Taylor coefficients of
each \(U_{y_j}\theta_0\) in \(\overline{\mathbb Q}\).
All its derivatives at zero are algebraic as well.
Each coordinate linear form \(z\mapsto(gz)_j\) has algebraic
coefficients by VI1. Thus every factor in (VI8.1), after this
substitution, is a literal \(\texttt{Gen.lin}\) generator for
the algebraic-coordinate directions.

The remaining coefficients in (VI7.1)--(VI7.2) are
\(\kappa\in\overline{\mathbb Q}^{\times}\), signs
\(\alpha(y)\), and \(r^{-1/2}\in\overline{\mathbb Q}^{\times}\).
The quotient is finite, and changing its representatives does
not change the function by VI4. The \(\texttt{Gen}\) rules for
constants, finite sums and products now prove
\[
 \texttt{Gen}\left(z\longmapsto\prod_{a=0}^3 f(z_a)\right).
 \tag{VI8.2}
\]
The function \(f\) is entire by VI7. Hence (VI8.2) is exactly
the source \(\texttt{Root4}(f)\), proving
\(\texttt{thetaVector\_fourthPower\_algebraicGerm}\).
There is no need to infer algebraic coefficients for an
unspecified complex-linear map from rational spanning alone:
VI1 supplied the actual algebraic matrix.

\paragraph{VI9: the same identity handles every rational Heisenberg translate.}
The preceding argument can be strengthened in a way useful for
the later indicator calculation. Let \(x\in\mathbb Q\mathcal M\)
and set \(f_x=U_xf\). Put \(x^{(4)}=(x,x,x,x)\) and
\(v=gx^{(4)}\). Some \(m>0\) has \(mv\in\mathcal L\).
Applying \(U_{x^{(4)}}\) to (VI7.2), using isometry of \(g\)
and the first identity in (VI4.2), gives
\[
 f_x^{\otimes4}(z)=
 \frac{\kappa}{\sqrt r}
 \sum_{[y]\in\mathcal L/(\mathcal L\cap\mathcal S)}
  \alpha(y)e^{-\pi iE_H(v,y)}
       \prod_j (U_{v_j+y_j}\theta_0)((gz)_j).
 \tag{VI9.1}
\]
The sign and the factor \(\pi\) in the phase are those of
the actual Heisenberg multiplication formula.
Since \(\mathcal L\) is integral,
\(mE_H(v,y)=E_H(mv,y)\in\mathbb Z\). Thus every displayed
phase is a root of unity and is algebraic. By (VI1.3),
every \(v_j+y_j\) is a rational-lattice point of \(\Lambda\).
The same TA10 argument as in VI8 proves
\(\texttt{Root4}(f_x)\).

For any multi-derivative \(D^\nu\) in algebraic coordinate
directions, differentiate (VI9.1) by that same operator
separately in each of the four blocks and then set all
variables to zero. TL18's product and chain rules show that
the resulting number is algebraic. Its left side is
\((D^\nu f_x(0))^4\), so \(D^\nu f_x(0)\) is algebraic.
This remains valid when \(f_x(0)=0\); no value of \(f_x\)
has been divided out. We have proved algebraicity of all
Taylor coefficients of every rational translate of this
actual vacuum, without invoking a theorem about arbitrary
theta functions with algebraic coefficients.

\paragraph{VI10: rectangular vacua and the actual moved functions.}
For arbitrary \(N,N'\), the standard matrix lattice satisfies
the exact integral equality
\[
 \mathcal A(\mathcal O_E^N,\mathcal O_E^{N'})
       =\mathcal M^{\,N'\times N}.
 \tag{VI10.1}
\]
An integral homomorphism between the standard modules is a
matrix over \(\mathcal O_E\). Multiplying by \(A\) gives
matrices with entries in \(A\); conversely every such
matrix is the sum of its individual matrix entries.
Applying the defining period generators proves (VI10.1)
in both directions.

The rectangular vacuum is exactly
\[
 f_{N,N'}(z)=\mathscr B.\texttt{thZero}(N,N')(z)
              =\prod_{i,j} f((z_{\tau ij})_\tau).
 \tag{VI10.2}
\]
The Hermitian form and its multiplier split over these
columns, so VI7 proves that it is entire, satisfies the
literal theta condition, and is nonzero. The one-copy
spanning assertion of VI7, applied successively to the
columns, proves that it spans the full standard theta line.
This also covers an empty matrix index set: the product is
the constant \(1\) on the zero-dimensional space.

If \(X\in\mathbb Q\mathcal A(\mathcal O_E^N,\mathcal O_E^{N'})\),
the Heisenberg operator on (VI10.2) is the product of the
operators on its columns. Every column belongs to
\(\mathbb Q\mathcal M\). VI9 therefore proves that all Taylor
coefficients of \(U_Xf_{N,N'}\) are algebraic. At \(X=0\) this
is the literal \(\texttt{thZero\_coefIn}\); division by the
multi-index factorials preserves algebraicity.

For any actual self-dual pair \(L,L'\), VI2 gives finite
commensurability with the standard period lattice and the
same multiplier. Both actual moving operators defining
\(\texttt{thMoved}\) are consequently the positive finite
sums (VI4.3). They preserve the theta condition and do not
kill the vacuum, by VI5. All their translating vectors lie
in the rational span of the standard period lattice.
Their coefficients are signs times inverse square roots
of positive integers.

This finite presentation is stable under a further rational
translation. For \(X,Y\) in the rational span, choose
integers \(m,n>0\) with \(mX,nY\) in the standard integral
period lattice. Then \(mnE_H(X,Y)\in\mathbb Z\), so the
phase in \(U_XU_Y=e^{-\pi iE_H(X,Y)}U_{X+Y}\) is algebraic.
Finite sums and VI9--VI10 now prove the actual assertion
\(\texttt{ta\_thMoved}\): the moved function is entire and
every rational translate has algebraic Taylor coefficients.
All sums here are genuinely finite for these lattices;
no support assertion is inferred from a totalized sum.

\paragraph{VI11: normalized indicator jets and smoothness on the same ambient data.}
Retain the literal seam \(s\), fields \(X\), character \(\chi\),
and constants \(c_D\) of \(\texttt{ambOf}(s,\mathscr B,\chi,c_D)\).
Assume exactly the source's \(X.\texttt{Ok}\), \(s.\texttt{Ok}\),
\(\texttt{Crv}\), \(\texttt{IsChi}\), and the stated even ranks.
The independent finite-presentation assertion CA2 gives, for
each allowed actual indicator,
\[
 \chi_F^{-1}\,\texttt{thInd}(z)
       =\sum_{\ell\in S} r_\ell\,
           U_{y_\ell}\bigl(\mathscr B.\texttt{thMoved}(L,L')\bigr)
                    (\operatorname{reindex}_s z),
 \qquad r_\ell\in\overline{\mathbb Q},
 \quad y_\ell\in\mathbb Q\mathcal A(L,L').
 \tag{VI11.1}
\]
CA2 proves this formula from the actual finite indicator;
it does not assume algebraic jets. VI2 identifies the
rational span on the right with that in VI10. Hence VI10
proves algebraicity of every coefficient of the right-hand
side. The seam reindexing has entries \(0,1\), and its
multi-derivative dictionary only introduces the explicit
integer factorials. Thus (VI11.1) proves the literal
\(\texttt{JetAlg}\) for every finite prime set and every
depth in the source quantifier order.

The ambient entire function is a constant multiple of the
same moved theta function after this coordinate reindexing.
VI10 proves its complex smoothness of every order, giving
the literal \(\texttt{ESmooth}\). Multiplication by a complex
constant preserves smoothness; no algebraicity of the
unnormalized factor \(\chi_F\) is silently assumed.

Consequently \(\texttt{esmooth\_jetAlg}\) is supplied by the
actual vacuum construction, the one-variable PA/TA input,
and CA2's independently proved finite identity. The dependency
direction is VI1--VI10 and CA2 into VI11; none of those
primitive assertions uses the final \(\texttt{FinAlg}\) or
\(\texttt{EvLg}\) construction. The arithmetic host and the
named classical one-variable analytic facts remain explicit;
no full RH or GRH conclusion is asserted here.
